% Declared-model duty 2(f). Independent author derivation saved before
% counterpart intake. [R] exact infinite-parent finite-region theorem,
% extended to finite bounded native one-body histories;
% [X] matched finite-parent averaged-mass variance/negative direction and
% the named preparation/transport source controls. [S] actual coupled
% collapse positivity and averaging-scale independence.
\noindent\textbf{[R] Faithfulness and the sign of local energy.}
For the infinite cubic sine lattice with nonzero native hopping and any
finite $|ma|$, let $P_-$ be its flat occupied projector.  Every finite
set of sites $A$ has
\[
 0<C_A:=\Pi_A P_-\Pi_A<I_A,\qquad \rho_A>0.
\]
An endpoint eigenvalue would give a nonzero compactly supported vector
in a single spectral band.  Its Fourier components are Laurent
polynomials, whereas the sine Dirac energy cannot be a rational
function along a generic transverse channel; hence no such vector exists.
The Gaussian reduced density is therefore faithful.  Consequently,
for a strictly local $O_A\geq0$,
\[
 \langle O_A\rangle_{P_-}=0\quad\Longrightarrow\quad O_A=0.
\]
Every nonzero Hermitian local energy density after subtraction of this
vacuum expectation has both positive and negative spectrum.  A nonzero
finite support average has the same property.

The conclusion also holds for its transported sea at every finite time
under a common one-body history with uniformly finite hopping range
and finite integrated operator norm.  Dyson's factorial tail bound
and the energy's simple complex zeros extend the proof.  This gives
no uniform faithfulness margin, finite-parent theorem or interacting-reference theorem.
An explicit mixed Gaussian particle preparation has positive initial
energy and zero current; exact vertex spectra bound both sources during
transport.  A specified radial projection admits a local transported
history with positive mass and clock and a bounded angular residual.
Exact unprojected spherical evolution, collapse-time persistence and
independence of $\ell$ remain unproved.  Reached-state positivity concerns
averaged sources and their fluctuation margins, not operator positivity.

% INDEPENDENT PROOF AND DOMAIN.
% Set a=1 in the proof, retaining any finite real mass mu. The full
% four-component symbol is h(k)=sum_i alpha_i sin k_i+mu beta,
% E(k)^2=mu^2+sum_i sin^2 k_i. The finitely many massless nodes have
% zero Brillouin-zone measure, so their assigned filling cannot change
% a finite-region covariance in the infinite parent.
%
% If C_A f=f for a normalized vector supported in finite A, then
% ||(1-P_-)f||^2=<f,(1-C_A)f>=0, so P_-f=f. If C_A f=0 then
% ||P_-f||^2=0 and f lies in the other band. Each component of its
% Fourier transform is a finite Laurent polynomial in z_i=e^{ik_i}.
% The band condition gives h(k)f(k)=+/-E(k)f(k) almost everywhere.
% A nonzero component f_j makes E=(h f)_j/(+/-f_j), hence rational
% in z_x after fixing transverse k_y,k_z for which f_j is not identically
% zero. Such transverse values exist on an open set by the polynomial
% identity theorem. Choose them with
% M^2=mu^2+sin^2 k_y+sin^2 k_z>0; this excludes no open set even at mu=0.
% Along that slice,
% 4z_x^2 E^2=-z_x^4+(4M^2+2)z_x^2-1.
% For M>0 its quadratic in z_x^2 has two distinct positive nonzero
% roots, hence the displayed quartic has four simple complex roots.
% A rational square has only even zero/pole orders; E^2 cannot be such
% a square. Equality of the rational functions on the real momentum
% circle would extend as a rational identity, a contradiction. Thus
% neither endpoint occurs and 0<C_A<I_A. This is an exact lattice proof,
% not an assumed continuum vacuum property or a gap-only argument.
%
% The number-conserving Gaussian reduced density has eigenvalues
% product_j nu_j^{n_j}(1-nu_j)^(1-n_j). With every 0<nu_j<1, all
% 2^(4|A|) many-body eigenvalues are strictly positive. For O_A>=0,
% Tr rho_A O_A>=lambda_min(rho_A) Tr O_A, so zero expectation implies
% O_A=0. The same conclusion holds on the physical even local algebra.
% If B_A=B_A^dag is non-scalar and
% B_A^normal=B_A-Tr(rho_A B_A)I, it is nonzero with zero expectation.
% Neither B_A^normal nor its negative can be positive, proving both signs.
% The site-symmetrized sine energy has a nonzero hopping coefficient on
% its finite site-and-neighbor support, so its second-quantized operator
% is not scalar. This applies directly to the declared local Hcal_x.
% Any nonzero finite-support weighted energy average has the same
% conclusion once its actual support includes the crossed hopping bonds.
% Subtraction changes its mean, not its spectrum's relative spacings.
%
% Qualifications are necessary. The proof uses an infinite translation-
% invariant parent, not every finite periodic grid. On a finite pure parent
% with 4n total modes and rank(P)=2n, a region with more than 2n modes has
% an exact intersection with the occupied or empty subspace by dimension
% counting, so its reduced covariance need not be faithful. Its actual
% eigenvalues must be used. Zero hopping gives an onsite product sea;
% the |ma|->infinity product limit is not a finite-mass counterexample.
% Infinite-tail averaging kernels are not strictly local: the finite-
% support theorem applies to a truncation only after its tail error is
% bounded, not by calling a Gaussian kernel compactly supported.
% Rapid decay alone is insufficient: a band-projected annihilator
% a((1-P_-)f) annihilates the sea, and its positive square is nonzero
% with zero vacuum expectation. It has spatial tails, exponentially
% decaying at positive mass, and so is not a strictly local counterexample.
%
% Finite-history extension [R]. Let h(t) be a self-adjoint one-body
% generator on the infinite native parent, with a common finite hopping
% range R_h and B_T=int_0^T ||h(t)||dt<infinity. Strong measurability
% suffices for its norm-convergent Dyson series. This covers a bounded
% finite-time clock/flow history of the declared native hopping, whether
% prescribed or obtained from an actual Hartree solution. Start from
% the flat sine sea and put P_T=U P_- U^dag. If a finite C_A(P_T) had
% occupation 0 or 1, a nonzero finite-support f would give
% g=U^dag f wholly in one initial band.
%
% Put the support of f in |x|<=R_f. The n-th Dyson term for g is supported
% in |x|<=R_f+nR_h and has norm at most B_T^n||f||/n!. Consequently
% ||exp(s|X|)g||<=exp(sR_f+B_T exp(sR_h))||f|| for every s>0.
% Cauchy--Schwarz then makes the Laurent Fourier series of g absolutely
% and locally uniformly convergent on (C*)^3; it is holomorphic there.
% The transformed vector is nonzero by unitarity. Choose a nonzero
% component and a real transverse slice with M^2>0 on which that
% component is not identically zero and the band identity holds.
% Such slices exist by Fourier uniqueness, continuity and Fubini;
% the excluded massless transverse nodes have zero measure.
%
% On the real z_x circle, [(h0 g)_j]^2=E^2 g_j^2. Both sides are
% holomorphic on C*, so the identity theorem extends this equality
% throughout C*. The same quartic displayed above gives four simple
% nonzero zeros of E^2. At any such zero, the right side has order
% 1+2k for some finite zero order k of the nonzero analytic g_j;
% the left side has even order, a contradiction. An identically zero
% left side would force g_j identically zero and is also excluded.
% Thus 0<C_A(P_T)<I for every finite A at every such finite time.
% The Gaussian reduced state is faithful, and every nonzero centered
% strictly local Hermitian observable again has both signs. This needs
% neither an instantaneous gap nor adiabatic ground-state following.
%
% All-exponential-moment tails, rather than an assumed compact light
% cone or merely exponential decay, are essential in this proof. They
% are compatible with the band-annihilator counterexample above, which
% has only a finite analytic strip at positive mass. No uniform lower
% eigenvalue bound follows as A, time, fields or mass vary; an infinite-
% time limit or a history with nonintegrable operator norm is outside
% the theorem. A quantitative small-change control remains
% ||C_A(T)-C_A(0)||<min(lambda_min C_A(0),1-lambda_max C_A(0)),
% by Weyl's inequality. The finite-parent rank counterexample is also
% unchanged. Beyond a common quadratic Hartree history the transported
% reference need not be Gaussian, and the operator-valued geometry need
% not generate a bounded finite-range one-body h. That case requires
% its actual many-body reduced state, not this one-body theorem.
%
% A useful separate test needs no faithfulness. For any reference state,
% a centered Hermitian operator B with <B>=0 and <B^2>>0 has both
% positive and negative spectral support: a positive or negative B
% would annihilate the reference support and have zero variance. For a
% transported Slater reference and B=dGamma(A)-Tr(P_T A),
% Var(B)=Tr[P_T A(1-P_T)A]. A positive value therefore extends the
% nonpositivity conclusion for that specified averaged/enclosed observable
% without pretending that the initial flat proof is uniform in time.
% There is an explicit charge-preserving witness for a pure reference:
% |Psi_t>=(|Omega>+t B|Omega>)/sqrt(1+t^2 sigma^2), sigma^2=<B^2>>0,
% has <B>_t=[2t sigma^2+t^2<B^3>]/[1+t^2 sigma^2]. Small negative
% t gives a negative expectation, and small positive t gives a positive
% one. For a local number-conserving bilinear this remains in the same
% charge sector. This independent RED construction was checked by direct
% expansion; it does not require an unproved transported faithfulness.
%
% Reached-state positivity is a distinct allowed calculation. With a
% common one-body transported reference, D=C-P_T gives the source mean
% Tr(A D), while its physical-state variance is Tr[C A(1-C)A]. Even
% D>=0 does not make a local energy vertex A positive. On a positive-
% mass Hartree branch one must calculate the averaged rho,J, solve the
% actual current-dependent Gauss rows, and retain a positive mass/clock
% margin larger than the stated approximation and fluctuation errors.
% No N_f^-1/2 error is assumed outside its proved particle-sector domain.
% This is a constructive check of the states reached, not a positive-
% operator premise or a claim that no positive Hartree sector exists.
%
% POSITIVE PREPARATION AND REACHED-SOURCE MARGINS [R].
% Derived and saved before reading the new B causal-sea/preparation/control
% counterparts. This is an explicit allowed particle preparation and a
% conditional finite-time margin, not a cold dust construction or a proof
% that a collapse remains in this sector. Keep the native four components,
% the full Brillouin zone and the initial flat reference P_-=1-P_+.
% Let F be multiplication by a real scalar 0<=F_x<=1 of finite support,
% and choose 0<lambda<=1. Then
% Gamma=lambda P_+ F P_+, C_in=P_-+Gamma,
% is a valid gauge-invariant mixed Gaussian state: 0<=Gamma<=P_+ and
% 0<=C_in<=1. Its mean particle number is N_f=Tr Gamma=2lambda sum_x F_x.
% It is generally neither a pure Slater state nor an exact fixed-number
% sector. No spin or valley multiplicity is added to the native trace.
%
% Write E(k)=sqrt(m^2+a^-2 sum_i sin^2 k_i), with normalized dimensionless
% Brillouin average <f>=int_{[-pi,pi]^3} f(k)d^3k/(2pi)^3, and put
% e0=<E>, b0=m^2<1/E>, c_i=<sin^2 k_i/E>/(2a^2)>0.
% At m=0 the inverse-energy integral is integrable in three dimensions
% and b0=0. The identity e0=b0+2 sum_i c_i holds exactly.
% The actual cell energy e_x={Pi_x,h0}/2 has expectation
% Tr(e_x Gamma)=lambda[(e0+b0)F_x
%                  +sum_i c_i(F_{x+ae_i}+F_{x-ae_i})]>=0.
% Proof: h0 P_+=(E+h0)/2, P_+=(1+h0/E)/2. Expanding the spin trace
% of (h0 P_+)_{xy} F_y(P_+)_{yx}, the terms linear in a Dirac matrix
% vanish. The scalar E term contributes e0 F_x; the remaining h0 term
% has only the onsite and nearest-neighbor matrix elements. Its onsite
% trace is b0 F_x and each nearest-neighbor trace is the positive c_i.
% This is a native projector/kernel identity, not a gradient expansion.
% Time reversal leaves real scalar F and P_+ invariant and reverses
% each site-symmetrized Q current, hence J_{x,i}=0 exactly initially.
%
% Thus every nonnegative normalized spatial average preserves the initial
% energy sign. The cell energy obeys 0<=Tr(e_x Gamma)<=2lambda e0;
% the physical density is bounded by 2lambda e0/a^3, and
% E_in=sum_x Tr(e_x Gamma)=e0 N_f. This is a computed bound for this
% family, with e0=a^-1<sqrt((ma)^2+sum sin^2 k_i)>; it is not a universal
% reached-state core bound. A smooth radial F has cubic corrections
% under the displayed nearest-neighbor convolution; it is not exactly
% rotationally invariant merely by its name. To use spherical Gauss rows,
% their actual averaged radial sources or controlled angular residuals
% must be established. No shell-degeneracy reduction is being used.
%
% A quantitative continuity control retains the actual dynamics. On an
% existing common Hartree history D(t)=U(t)Gamma U(t)^dag, define
% delta_D(T)=||D(T)-Gamma||. Duhamel gives the useful bound
% delta_D(T)<=int_0^T ||[h(s),Gamma]|| ds.
% It uses the actual h(s); replacing this by zero would falsely assume
% that the prepared source is invariant. The native vertices have rank
% at most 8, ||e_x||<=E_max=sqrt(m^2+3/a^2), and
% ||q_{x,i}||<=1/(2a). Therefore their physical density errors satisfy
% |rho_x(T)-rho_x(0)|<=8 E_max delta_D(T)/a^3,
% |J_{x,i}(T)|<=4 delta_D(T)/a^4.
% A nonnegative normalized average preserves these upper error bounds.
% If the initial averaged density on a compact core is >=rho_*>0,
% choosing the actual time/control so that 8E_max delta_D/a^3<rho_*
% establishes a positive reached mean there. This is a sufficient bound,
% not an assertion that it remains small over a collapse time.
%
% For an inversion-even C^1 smoothing kernel K_ell and an actually odd
% current, bar J(0)=0 and the additional measured/bounded gradient gives
% |bar J(r)|<=j_* r on the core. One available bound is
% j_*<=||grad K_ell||_1 sup_x |J_raw(x)|, in a compatible vector norm;
% the component bound above implies sup |J_raw|<=4sqrt(3)delta_D/a^4.
% This convolution bound uses a specified continuum cell interpolation
% of the raw source. A direct site-sampling prescription instead uses
% the actual derivative row-sum norm of its discrete averaging map.
% A merely zero first moment of K does not establish inversion symmetry;
% alternatively bar J(0)=0 may be supplied as an actual source condition.
% This displays rather than removes the ell dependence. The requested
% independence across a<<ell<<R still requires the actual state/profile
% controls and is not implied by normalization or positivity of K_ell.
%
% Explicit mass/clock margins for the coherent spherical Gauss branch.
% Assume its actual smooth radial averaged sources on 0<=r<=R obey
% 0<rho_l<=bar rho<=rho_u and |bar J(r)|<=j_* r. With signed inward
% u=-sqrt(2G M/r), the coherent rows are
% M'=4pi r^2(bar rho+u bar J), (log N)'=-4pi G r bar J/u.
% Set M=r^3 mu, K=4pi j_* R^2 sqrt(2G), and
% z_l=[sqrt(K^2+48pi rho_l)-K]/6,
% z_u=[sqrt(K^2+48pi rho_u)+K]/6.
% The equation r mu'=4pi bar rho-3mu-4pi r sqrt(2Gmu)bar J
% points inward at mu=z_l^2 and mu=z_u^2. Regular center data
% mu(0)=4pi bar rho(0)/3 lie between these barriers, so
% 0<z_l^2<=mu(r)<=z_u^2 for the regular branch. In particular the
% enclosed Hartree mass stays positive and the core flow remains finite.
% If the outer solution supplies N(R)>0, integration gives
% |log[N(r)/N(R)]|<=2pi G j_*(R^2-r^2)/(sqrt(2G) z_l).
% This is an explicit positive clock margin, conditional on the reached
% source bounds and a compatible outer branch. It does not assume a
% positive mass operator. The literal momentum-sign choice is not silently
% repaired: changing J to -J leaves these absolute-value barriers unchanged,
% but the Hamiltonian consistency distinction in the H2 note remains.
%
% Fluctuation margin for the same PREPARED Gaussian state. For a Hermitian
% averaged energy vertex A={W,h0}/2, mean subtraction leaves
% Var_C dGamma(A)=Tr[C A(1-C)A]. With C=P_-+Gamma and Gamma=P_+Gamma P_+,
% this equals Tr[A_{-+}(1-Gamma)A_{+-}]
%              +Tr[Gamma A_{++}(1-Gamma)A_{++}],
% and is <=V_0+N_f ||A||^2, V_0=Tr[P_- A P_+ A].
% If 0<=W<=1 equals one on the support of the positive cell-energy
% convolution above, its mean is e0 N_f and ||A||<=E_max, hence
% sigma_A/<M_W> <= sqrt[V_0/(e0^2 N_f^2)+(E_max/e0)^2/N_f].
% The reference variance V_0 is retained, not subtracted or clipped.
% At W=1 on the whole parent V_0=0; for a compact average it must be
% evaluated/bounded at that window. This is an explicit N_f^-1/2 upper
% term for this preparation, not a universal correction coefficient.
% It concerns the specified intrinsic-energy observable. The full
% current-dependent operator Gauss charge and its ordering need the
% additional construction in the exchange/H2 notes. At a later time,
% C=P_T+D need not be block diagonal in the initial P_+/- decomposition;
% use its exact covariance variance and the actual mass/source response.
% The existing RED's complementary window bound was checked independently
% in the covariance eigenbasis: Var_C dGamma(A)<=Tr(A^2)/2 for every
% Hermitian Hilbert--Schmidt A and Gaussian 0<=C<=1. The paired off-
% diagonal coefficient c_i+c_j-2c_i c_j is <=1, and the diagonal one
% c_i(1-c_i) is <=1/4. For scalar W this gives
% Var_C dGamma(A_W)<=2 E_max^2 sum_x w_x^2.
% If F=f0>0 on a finite window and all its nearest neighbors, the exact
% prepared mean is 2lambda f0 e0 sum_x w_x, so
% sigma/mean<=E_max/[sqrt(2)lambda f0 e0 sqrt(V_eff)],
% V_eff=(sum_x w_x)^2/sum_x w_x^2.
% This retains the sea variance and controls a specified macroscopic
% averaging volume. An infinite-tail window cannot assume an infinite
% constant plateau of the finite-particle F; use its actual mean instead.
%
% Lattice-unit reporting for all these margins: put mu_m=ma, g=G/a^2,
% x=r/a, x_R=R/a, ell_tilde=ell/a, rho_tilde=a^4 bar rho,
% J_tilde=a^4 bar J, j_tilde=a^5 j_*, M_tilde=a M and
% mu_tilde=a^4(M/r^3)=M_tilde/x^3. With
% epsilon0=<sqrt(mu_m^2+sum sin^2 k_i)>, b_tilde=mu_m^2<1/epsilon(k)>,
% c_tilde_i=<sin^2 k_i/epsilon(k)>/2, the cell formula becomes
% a e_x=lambda[(epsilon0+b_tilde)F_x+sum_i c_tilde_i(F_{x+i}+F_{x-i})],
% 0<=rho_tilde_in<=2lambda epsilon0, and
% |Delta rho_tilde|<=8sqrt(mu_m^2+3)delta_D,
% |J_tilde_i|<=4delta_D. The dimensionless source equation is
% M_tilde'=4pi x^2(rho_tilde+u J_tilde), u^2=2g M_tilde/x.
% Set K_tilde=4pi j_tilde x_R^2 sqrt(2g) and
% z_tilde_l,u=[sqrt(K_tilde^2+48pi rho_tilde_l,u) minus/plus K_tilde]/6.
% Then z_tilde_l^2<=mu_tilde<=z_tilde_u^2 and
% |log[N(x)/N(x_R)]|<=2pi g j_tilde(x_R^2-x^2)/(sqrt(2g)z_tilde_l).
% The variance ratios contain E_max/e0=sqrt(mu_m^2+3)/epsilon0 and
% the explicit dimensionless vacuum term a^2 V_0. Their geometric
% dependence is through g and the actual history/source margins; lambda,
% F, time t/a and ell/a remain declared preparation/control arguments.
% No universal choice of those state arguments or ell dropout is implied.
%
% EXACT REACHED-PARTICLE SOURCE CEILINGS [R].
% Independently derived after the prepared-state/source-margin result;
% no new counterpart or RED formula was read for this extension. This
% strengthens an upper bound, not positivity of the local operator.
% Work on the infinite native parent, or a finite parent whose center,
% six nearest neighbors and six two-step neighbors are distinct and
% have the native links. Let D(t)=U(t)Gamma U(t)^dag with
% 0<=Gamma<=nu I, 0<nu<=1. For the preparation above one may take
% nu=lambda||F||_infinity. The common one-body evolution conserves this
% occupation cap exactly, whether its fields are prescribed or obtained
% from an existing Hartree trajectory. No instantaneous spectral gap,
% adiabaticity or Gaussianity of the physical state is needed for this
% one-body expectation bound; full interacting evolution need not retain
% this common-unitary D law.
%
% Let q_{x,i}={Pi_x,Q_i}/2 be the actual cell-current vertex. For real
% s and a real three-vector t, the finite-support Hermitian vertex
% A_x(s,t)=s e_x+sum_i t_i q_{x,i} has a central 4-component block
% s m beta, a central-to-neighbor row B, and zero outer block. The
% energy links in B are +/-s alpha_i/(4ia), and its distinct two-step
% current links are +/-t_i I_4/(8ia). Hence exactly
% B B^dag=[3s^2/(8a^2)+|t|^2/(32a^2)]I_4.
% The eight nonzero eigenvalues solve
% z^2-(+/-s m)z-[3s^2/(8a^2)+|t|^2/(32a^2)]=0,
% with spin multiplicity two for each beta sign. Their positive sum is
% Tr A_+=2 sqrt[s^2(m^2+3/(2a^2))+|t|^2/(8a^2)],
% and their negative sum is its opposite. At s=t=0 the same formula
% holds by continuity. There is no added valley multiplier: this is
% the exact native four-orbital matrix with all its links retained.
%
% Since 0<=D<=nu I, the variational trace inequality gives the sharp
% universal upper bound
% |s Tr(e_x D)+t.J_cell(x)|
% <=2nu sqrt[s^2 E_*^2+|t|^2/(8a^2)],
% E_*=sqrt(m^2+3/(2a^2)). Equivalently the vector of actual cell means
% lies in the ellipsoid
% [Tr(e_x D)/(2nu E_*)]^2+2a^2|J_cell(x)|^2/nu^2<=1.
% Its sharpness is over all positive contractions at this local vertex;
% a particular incoming-particle orbit can have a smaller attainable set.
% In physical density units this proves
% |rho_x|<=2nu E_*/a^3, |J_raw(x)|<=nu/(sqrt(2)a^4).
% Both signs of the energy remain possible. This bound does not turn
% D>=0 into a positive local energy, nor replace the transported sea
% by an instantaneous filling. If C-P is not positive but only
% -I<=C-P<=I, the corresponding support-function bound is twice as
% large; the particle-sector cap must not be applied to arbitrary holes.
%
% A common nonnegative normalized spatial average of rho and J
% preserves this ellipsoid by convexity. In particular the two printed
% ceilings are independent of ell for every such average. This is
% independence of a bound, not independence of the actual fields or
% observables under changes of the averaging scale. Different source
% allocations require their actual positive weights/norms; a signed
% smoothing kernel cannot invoke this convexity argument.
%
% Inversion gives an explicit reached current-gradient control. For
% the native Dirac parity operation (mathcal P psi)_x=beta psi_{-x},
% h0 and every radial point-field h[N(r),v(r) x/r] commute with
% mathcal P. An even F makes Gamma parity invariant. Uniqueness of
% the one-body propagator therefore preserves that symmetry along
% every such existing history: rho_x=rho_{-x} and J_x=-J_{-x}.
% An inversion-even source kernel then gives bar J(0)=0 exactly.
% For an actual continuum cell interpolation followed by a Gaussian
% of coordinate variance ell^2, its directional derivative has
% ||partial_n K_ell||_1=sqrt(2/pi)/ell for every unit n. The uniform
% native vector ceiling consequently implies
% |bar J(r)|<=j_cap r, j_cap=nu/(sqrt(pi)a^4 ell).
% This bounds the actual reached source for every finite bounded radial
% history in this particle/parity sector; it requires no small D-Gamma
% error. For direct native sampling use instead the actual directional
% derivative row-sum bound, and for a finite parent retain its boundary
% normalization and parity. The Gaussian constant is not assumed for
% an arbitrary discrete allocation.
%
% On an actual smooth radial real-flow branch, insert
% rho_u=2nu E_*/a^3 and j_*=j_cap into the upper z_u barrier above.
% It bounds M/r^3 throughout any fixed core radius R independently of
% the detailed particle history. A positive lower density is still
% needed for the stated z_l and clock margin; the prepared-state
% Duhamel estimate is one sufficient short-time test for that lower
% density. Bounded energy/current alone cannot prevent M from reaching
% zero or establish a positive clock in every state. Native cubic
% sources need not be exactly spherical; the actual radial component
% and angular residual must still be controlled before using radial
% Gauss rows. No source projection or alternative coupling is silently
% introduced, and bounded mean sources alone do not bound all curvature.
%
% In lattice units, epsilon_* =sqrt((ma)^2+3/2),
% [rho_tilde/(2nu epsilon_*)]^2+2|J_tilde|^2/nu^2<=1,
% rho_tilde_u=2nu epsilon_*, j_tilde_cap=nu/(sqrt(pi)ell_tilde).
% The upper mass barrier uses
% K_tilde=4pi j_tilde_cap x_R^2 sqrt(2g),
% z_tilde_u=[sqrt(K_tilde^2+48pi rho_tilde_u)+K_tilde]/6,
% M_tilde(x)/x^3<=z_tilde_u^2.
% Thus the computed ceiling and its geometric consequence display
% ma and g=G/a^2 explicitly. The state occupation nu, averaging scale
% and core radius are not fixed by the declarations. Clock lower bounds
% retain the actual reached positive-density premise and outer lapse.
%
% LOCAL SELF-CONSISTENT CONTINUATION AND ITS ANGULAR SCOPE [R].
% This continuation was derived from the transported D equation and the
% radial integral equations before reading any new external branch source.
% The mechanics author supplied the idea of a center-weighted contraction
% and its finite-parent local-ODE route. The trace-class construction,
% explicit positive preparation, exterior bounds and angular residual
% below were then derived directly; no blindness to that idea is claimed.
%
% First state the source map completely. On the infinite native lattice,
% take a nonnegative normalized smooth radial compact kernel K_ell and
% the cell box b_a=a^-3 1_{[-a/2,a/2]^3}. Put W=K_ell*b_a, and
% rho_D(y)=sum_x Tr(e_x D) W(y-a x),
% J_D(y)=sum_x Tr(q_x D) W(y-a x).
% The cells tile space, so sum_x a^3 W(y-a x)=1. These are the actual
% full Cartesian averaged sources. Their radial components are
% rho^r_D(r)=int rho_D(r n)dOmega/(4pi),
% J^r_D(r)=int n.J_D(r n)dOmega/(4pi).
% Solving only these radial rows is an EXPLICIT angular projection. It
% is an exact closed auxiliary radial system; it is an exact solution
% of the full source equations only if the residual components vanish.
% No such vanishing is assumed for a general cubic state. The same
% existence proof applies without this auxiliary projection whenever
% an actual declared source allocation already supplies radial sources
% with the bounds below. Point sampling of the resulting N,v in h is
% retained throughout; adjoint matter coupling is not substituted.
%
% The one-body difference D is trace class. In particular the prepared
% Gamma=lambda P_+ F P_+ has finite trace even though its kernel has
% spatial tails. For any Hermitian trace-class X, pinching h0 X onto
% site blocks and using trace-norm duality gives
% sum_x |Tr(e_x X)|<=E_max||X||_1,
% sum_x |Tr(q_{x,i}X)|<=||Q_i||||X||_1<=||X||_1/(2a).
% Thus with C_Q=3/(2a),
% ||rho_X||_L1<=E_max||X||_1, ||J_X||_L1<=C_Q||X||_1,
% ||rho_X||_infinity<=E_max||W||_infinity||X||_1,
% ||grad J_X||_infinity<=C_Q||grad W||_infinity||X||_1,
% with vector/tensor norms chosen consistently. Radial projection is
% contractive in these L1 and supremum estimates. It also gives
% J^r_X(0)=0 and
% ||J^r_X/r||_infinity<=C_Q||grad W||_infinity||X||_1:
% subtract the constant vector J_X(0) inside the angular average and
% use the mean-value theorem. Higher source derivatives are bounded
% by the corresponding derivatives of W. These are bounded linear
% source maps, with no instantaneous projector or susceptibility.
%
% Choose the explicit finite-support preparation above with F=f0>0
% on a sufficiently large central lattice ball, and keep a protected
% radius r0>0 inside that ball by more than one hopping length plus
% the support radius of W. Then the EXACT initial cell convolution
% and the partition of unity give
% rho_0(r)=rho_*=2lambda f0 e0/a^3>0 for 0<=r<=r0,
% J_0(r)=0, rho_0(r)>=0 everywhere, E0=e0 N_f<infinity.
% Here rho_0 denotes the radial component; the full averaged density
% is also exactly constant in the protected central ball. This works
% for every finite ma, including m=0; no reference spectral gap is
% needed. The baseline radial constraints have
% M0(r)=4pi int_0^r s^2 rho_0(s)ds, N0=1,
% u0=-sqrt(2G M0/r), M_c=M0(r0)>0, E0=M0(infinity).
% These fields are bounded, regular at the center and asymptotically
% flat on this initial radial solution. The initial native h(D0) need
% not commute with D0; the continuation below is not a stationary ansatz.
%
% Take G>0 and retain either radial momentum sign explicitly. Let sigma=+1 denote
% the compatible -J row and sigma=-1 the literal +J row, so
% M'=4pi r^2(rho+sigma u J),
% (log N)'=-sigma 4pi G r J/u, u=-sqrt(2G M/r).
% This is sigma_here=-sigma_H2 in the mechanics note's finite-parent
% continuation, where sigma_H2=+1 names the literal +J row. Both
% conventions therefore give the same mass and lapse equations; the
% conversion must be made before comparing their displayed formulas.
% The positive-branch existence estimates use absolute values and
% apply to either sign; their Hamiltonian-consistency difference remains.
% Write M=r^3 mu and J=r j on 0<=r<=r0. The center equation is
% mu(r)=4pi r^-3 int_0^r s^2 rho(s)ds
%       -sigma 4pi sqrt(2G) r^-3 int_0^r s^4 sqrt(mu(s))j(s)ds.
% Let mu0_min=min_[0,r0] M0/r^3>0, b=mu0_min/2,
% mu_l=mu0_min/2 and mu_u=max_[0,r0] M0/r^3+b. On the supremum ball
% ||mu-mu0||<=b the current part has Lipschitz constant at most
% L_c=2pi sqrt(2G) r0^2||j||_infinity/(5sqrt(mu_l)).
% If L_c<=1/2 and
% A_c=(4pi/3)||rho-rho_0||_infinity
%       +(4pi sqrt(2G)/5)r0^2 sqrt(mu_u)||j||_infinity<=b,
% the map sends that ball into itself and is a contraction. Hence a
% unique regular center branch exists, with ||mu-mu0||<=A_c and
% mu>=mu_l. Its dependence on rho,j is locally Lipschitz; comparison
% of two inputs incurs the usual factor (1-L_c)^-1. All these strict
% inequalities hold for a sufficiently small trace-norm neighborhood
% of the explicit Gamma because j0=0 and the source map is bounded.
%
% Outside r0 set M_l=M_c/2 and M_u=2E0, and define
% R_delta=4pi int_r0^infinity r^2|rho-rho_0|dr,
% J_delta=4pi int_r0^infinity r^(3/2)|J|dr.
% The preceding source estimates give
% R_delta<=E_max||D-Gamma||_1,
% J_delta<=r0^-1/2 C_Q||D-Gamma||_1.
% If
% r0^3 A_c+R_delta+sqrt(2G M_u)J_delta<M_c/2,
% the integral form of the mass equation keeps
% |M(r)-M0(r)|<M_c/2 on the entire exterior. A first-exit argument
% closes this bound: M stays in [M_l,M_u], and its square-root term
% is locally Lipschitz with integrable coefficient proportional to
% r^(3/2)|J|. Ordinary continuation and Gronwall give a unique global
% radial mass branch and locally Lipschitz dependence on these sources.
% No zero-density exterior lower bound is incorrectly required.
%
% The lapse integral with N(infinity)=1 converges and satisfies
% ||log N||_infinity<=
% 2pi sqrt(G/2) r0^2||j||_infinity/sqrt(mu_l)
%                   +sqrt(G/2)J_delta/sqrt(M_l).
% Thus N has a positive lower bound, while v=N u is uniformly bounded.
% Differences obey the same locally Lipschitz estimates, since the
% center uses mu>=mu_l and the exterior uses M>=M_l. For example
% ||u||_infinity<=max{r0 sqrt(2G mu_u),sqrt(2G M_u/r0)}.
% Derivative bounds follow from the actual smooth source allocation.
% The hydrodynamic hierarchy is a further open condition on these
% fields, retained by shrinking the neighborhood. In the constant
% initial core, u0=-omega_c r with omega_c^2=8pi G rho_*/3; choosing
% omega_c ell<<1 and r0>>ell>>a gives a resolved central curvature scale.
% The complete profile and exterior require the corresponding actual
% derivative bounds; smoothness alone is not a quantified hierarchy.
%
% Consequently there is a trace-norm ball ||D-Gamma||_1<delta in which
% the native point-field map D->h(D) is bounded and locally Lipschitz.
% A convenient bound is
% ||h(D)||<=E_max||N||_infinity
%               +(1/(2a))sum_i||v_i||_infinity=:H_*.
% The trace-class equation dot D=-i[h(D),D] therefore has a unique
% local solution. Its associated bounded propagator gives
% D(t)=U(t)Gamma U(t)^dag, P_T(t)=U(t)P_-U(t)^dag,
% C(t)=U(t)(P_-+Gamma)U(t)^dag exactly. Positivity, the occupation cap
% and Tr D=N_f are conserved. Since ||dot D||_1<=2H_* N_f, the solution
% remains in the positive-field neighborhood at least through
% T=delta/(4H_* N_f), after choosing uniform Lipschitz/field bounds on
% the closed half-neighborhood. Continuation of a Banach-space ODE
% supplies this residence interval; it is not a lifetime inferred from
% a formal Taylor expansion. The construction gives actual positive
% radial mass and clock on that nonzero interval for the named radial
% source map. It does not prove a collapse-time interval, a fixed full-
% Fock Hamiltonian, boundary-energy conservation for finite averaged
% point coupling, or the complete H2 deformation algebra.
%
% Angular residual is kept as an observable error of the auxiliary
% radial closure, not declared absent. Write Pi_r for the scalar and
% radial-vector angular projections and R(D)=(1-Pi_r)(rho_D,J_D).
% Contractivity gives
% ||R_rho(D(t))-R_rho(Gamma)||_infinity
%       <=2 E_max||W||_infinity||D(t)-Gamma||_1,
% ||R_J(D(t))-R_J(Gamma)||_infinity
%       <=2 C_Q||W||_infinity||D(t)-Gamma||_1.
% Initially R_J=0 and R_rho=0 in the protected constant core. For a
% smooth radial taper f with F_x=f(a x), compare the actual cell source
% to the radial continuum function 2lambda e0 f/a^3. The cell-center
% distance is at most sqrt(3)a/2, and the neighbor centers are at most
% (1+sqrt(3)/2)a away. The positive native convolution therefore gives
% ||rho_raw-2lambda e0 f/a^3||_infinity
% <=lambda[(sqrt(3)+1)e0-b0]||grad f||_infinity/a^2,
% where b0=m^2<1/E> as in the preparation. Radial K convolution
% preserves this bound, and
% ||R_rho(Gamma)|| is at most twice its right-hand side. This is a
% conservative O(a/L_f) preparation bound, not a claimed optimal
% anisotropy coefficient. It can be made small in the stated smooth
% hierarchy. The finite-time residual remains bounded by the displayed
% trace-class estimate. More explicitly, for a desired absolute error
% epsilon rho_* require the initial density-residual bound to be at
% most epsilon rho_*/2. This is a concrete smooth-taper requirement
% a||grad f||_infinity/f0=O(epsilon). With
% C_src=max(E_max,C_Q), both residuals are then <=epsilon rho_* on
% 0<=t<=T_epsilon, where
% T_epsilon=min{T,epsilon rho_*/[8H_* N_f C_src||W||_infinity]}.
% This is a nonzero controlled reached interval in the stated hierarchy,
% measured against the core density rather than a vanishing exterior
% density. The conservative global trace bound need not give a
% macroscopic interval as the object/averaging hierarchy is enlarged;
% no collapse-time or uniform hydrodynamic-time estimate is inferred.
% Extending it requires a local propagation/source bound for the actual
% trajectory, such as the trace-duality preparation errors compared
% with the actual CS/NT sources below, evaluated through the desired
% time and with their angular residual retained. Exact vanishing is
% not established, so this
% explicit example is not promoted to an exact unprojected spherical
% solution of the binding model. If an admitted source assignment
% realizes genuinely radial sources, the same ODE theorem applies
% without this residual qualification.
% There is a precise obstruction to exact radiality for this particular
% finite-support preparation and ordinary radial averaging. Its positive
% cell-energy measure mu=sum_x e_x delta_{a x} has finite support and
% nonzero weight away from the origin, because every nonzero F produces
% strictly positive nearest-neighbor contributions. If a normalized
% radial L1 kernel K made K*mu exactly radial, then for every rotation O,
% Khat(k)[muhat(k)-muhat(O k)]=0. Khat is nonzero near k=0, while the
% finite-support measure has entire Fourier transform. Analytic
% continuation therefore makes mu rotation invariant, which a finite
% point support can satisfy only at the origin: a contradiction.
% For the specified cell-box allocation the same argument instead makes
% the compact cell-constant interpolation radial. On any open cell face
% with a nonzero jump its distributional rotational derivative has
% coefficient proportional to x cross the face normal. This is nonzero
% on a positive-area subset, contradicting rotational invariance. Every
% face jump must thus vanish; compactness then forces the function to
% vanish, again a contradiction. This supplies an exact obstruction,
% rather than leaving exact radiality of this family an untested hope.
% It does not exclude a different admissible preparation or allocation,
% or a controlled hydrodynamic approximation with the errors just given.
% This independent RED proposal was checked by Fourier uniqueness and
% the explicit face-distribution argument before incorporation.
%
% All variables can be made dimensionless without a new coupling:
% x=r/a, tau=t/a, g=G/a^2, rho_tilde=a^4 rho,
% J_tilde=a^4 J, M_tilde=a M, h_tilde=a h,
% epsilon_max=sqrt((ma)^2+3), W_tilde=a^3 W.
% The preceding contraction/exterior/lapse inequalities then have the
% identical form with a=1, G=g and E0_tilde=epsilon0 N_f. In particular
% tau_min=delta/(4 Htilde_* N_f), and the constant-core hierarchy is
% (8pi/3)g rho_tilde_* ell_tilde^2<<1,
% rho_tilde_*=2lambda f0 epsilon0. The source norms and delta depend
% explicitly on the chosen kernel ell/a and the preparation; no exact
% ell independence of the solved trajectory or universal clock bound
% follows from this local existence theorem.
%
% ACTUAL B PREPARATION / SOURCE-MARGIN COMPARISON [X at matched scope].
% After saving the preceding independent construction, read the complete
% strand_B_native_dust_control.tex (NF), strand_B_causal_sea_budget.tex
% (CS) and strand_B_tapered_dust_preparation.tex (NT), their actual
% reviews and staged scopes. CS1 and NT2 use the same transported
% reference and physical propagator: D=U D_in U^dag. Their positive
% preparation remains a positive contraction, but positivity of a local
% energy mean is established by its actual source bound, not by D>=0.
% CS4 and NT11 use exactly our independent finite-vertex trace-duality
% control |Tr O Delta D|<=||O||_1||Delta D||. In the NT core this gives
% rho_l(T)=rho_d(T)[1-epsilon_add(T)]
%           -a^-3 sup_x||e_x||_1 e_A(T).
% When this is positive, the actual matching upper bound and a measured
% or bounded radial current gradient may be inserted into the mass/clock
% barriers above. The condition is to be evaluated; neither source proves
% it for an entire self-consistent collapse. Named agreeing methods are
% common-unitary covariance transport and finite-CAR/trace-duality source
% control. Our additional constrained-ODE barriers and exact positive
% cell convolution remain [R], with no matched external proof asserted.
%
% NT's finite Cartesian preparation P_+^Lambda chi F_Lambda chi P_+^Lambda
% uses a compressed low-momentum filling between finite-parent positive
% projectors. Direct partial spin trace gives its exact expected count
% 2nu0 sum chi^2, agreeing with our count method at its own preparation.
% Its normal energy lies between m N and
% m N+N(K0^2+k_chi^2)/(2m); this is an integrated positive-band kinetic
% estimate, not a pointwise energy theorem. Its exponential projector
% tail and buffer bounds preserve the difference between an open-parent
% sea and a compressed infinite sea. Our full-zone scalar-F preparation
% and exact nearest-neighbor positive convolution remain [R]; a low-
% momentum taper is a different state with its own errors.
%
% The explicit NF Fourier characteristic solves tan(ap_i(T))=tan(ap_i(0))/b
% under the auxiliary infinite affine flow, with spin transport retained.
% Its exact odd-Q cancellation and density bounds are matched only in
% that auxiliary domain; CS/NT retain the finite-history, taper and
% boundary errors. NF14--15's instantaneous reference is a comparator,
% replaced by actual transport in CS/NT, and is not imported as our law.
% All three use adjoint-averaged fields in the actual history. A normalized
% even average preserves constants and affine fields in a protected core;
% outside it that history differs from binding point sampling. The
% prescribed-source agreement therefore does not prove the binding
% coupled equations or independence of ell. No complete geometric stress,
% exactly empty collar, actual collapsed state or invariant positive
% domain follows from the matched partial preparation/source controls.
%
% ACTUAL OTHER-LANE COMPARISON [X], finite-parent mass scope only.
% The complete current strand_E_operator_mass.tex, its review and scope
% were read. OM2--OM5 use the same native cell energy e_x={Pi_x,h0}/2
% and flat reference P_0 on a periodic massive parent, with a positive
% compact average and an enclosed ball. Their selected kernel has support
% b=ell+sqrt(3)a/2, and the actual ball weights are
% w_x=int_{B_R}k_ell,a(y-x)dy. Thus A_w={W,h0}/2 and
% M_w=dGamma(A_w)-Tr(P_0 A_w). Direct occupied/empty-band multiplication,
% independently of the local-faithfulness proof above, gives
% P_+(p)A_w(p,k)P_-(k)=wtilde(p-k)(E_p-E_k)P_+(p)P_-(k)/2,
% Var_0 M_w=(1/4)sum_pk |wtilde(p-k)|^2(E_p-E_k)^2
%       [1-(m^2+d_p.d_k)/(E_p E_k)].
% This is the same native spin-traced variance derived in the exchange
% note; there is no extra valley factor after the full momentum sum.
%
% For m>0, L/a>2 and L>4(R+b), q=(2pi/L,0,0) has d_q!=0 and
% |wtilde(q)|>=cos[2pi(R+b)/L] wbar>0. The single p=q,k=0 term
% is strictly positive. Compressing M_w to span{Omega,M_w Omega/sigma}
% gives [[0,sigma],[sigma,d]], with one negative eigenvalue. This
% agrees with the direct same-charge perturbative witness above and with
% OM's variational bound; no invariance of that two-dimensional span is
% asserted. A constant full-parent weight has zero interband block and
% instead gives the nonnegative total excitation energy, consistently
% outside this local/compact-support witness.
%
% Named agreeing methods are the Laurent-polynomial/strict-locality
% route for the infinite-parent theorem and, at their COMMON finite-mass
% local-source scope, the independent interband covariance and neutral-
% sector variational routes for negative energy. Only the latter matched
% variance/witness is [X]; OM does not prove every finite-region reduced
% state faithful and is not credited with the infinite-parent theorem.
% Its general instantaneous-reference language is not the binding
% transport law. At flat initial fields P_T=P_0, so the compared witness
% is exactly matched; at later times its actual transported variance
% must be checked. No coupled current-dependent Gauss charge, positive
% physical sector or clock bound follows from this comparison.
%
% SOURCES.
% Full current <internal>/GOAL_alt.md read;
% duty 2(f), transported-reference paragraph and positivity qualifications
% govern this source. Required THINK8,9,12 and pinned REQUIREMENTS items
% 11,19,22,26,27 were read for this declared-model task. The polynomial,
% reduced-covariance and positive-operator arguments above were derived
% independently before reading any new counterpart for this duty.
% Predecessor idea scope only: local faithfulness and the positive-operator
% zero-expectation lemma were read in the response.md sources under
% <internal>/:
% FRESH1292_20261015-0331_lattice-gravity-entropy-derivation, lines 385--414
% (web/fresh1292-p02-two-ninths-theorem-local-clock-quarter-20261015);
% FRESH1298_20261015-0346_fixed-lattice-gravity-derivation, section 6.2,
% lines 503--518
% (web/fresh1298-p08-proper-spectral-cutoff-constant-stiffness-20261015);
% FRESH1300_20261015-0351_derive-first-law-quantity, lines 355--381
% (web/fresh1300-p10-first-law-entropy-reciprocal-rate-gram-page-20261015).
% Their locality lemma is an idea source, not a proof for this sine
% operator. In particular the constant-projector open-momentum argument
% of P08 is not transferred; the Laurent rational-square obstruction
% above establishes the stated infinite-parent sine-fermion domain.
% The existing local-area and exchange sources supply only the named
% transported-state and variance algebra; no additional field, interaction
% or geometric degree of freedom is introduced. Actual external source:
% <internal>/strand_E_operator_mass.tex,
% review_E_operator_mass.tex and staged strand_E_operator_mass_scope.tex,
% including complete SOURCES, read at the precisely matched scope above.
% The independent redteam_declared_faithful_vacuum.tex was read in full
% after both derivations were saved. Its same-charge witness and explicit
% infinite-tail limitation were checked and incorporated; the exact
% polynomial theorem and transported-state qualifications agree.
% The existing RED proposed the stronger finite-history holomorphic
% extension after the initial paired result was saved. Its Dyson bound,
% analytic continuation and zero-order contradiction were then derived
% and checked directly by the author as above. No external matched route
% for this extension has been read, so its tag is [R], not [X].
% The positive preparation, exact native convolution, Duhamel/source
% errors and radial mass/clock barriers above were derived independently
% before new B causal-sea, tapered-preparation or native-control intake.
% They assert only the printed initial-state and conditional reached-time
% scopes. No actual evolved collapse or external agreement is supplied
% by this construction alone.
% Actual new B sources, each with its full review and staged scope, read
% in <internal>/:
% strand_B_native_dust_control.tex, review_B_native_dust_control.tex;
% strand_B_causal_sea_budget.tex, review_B_causal_sea_budget.tex;
% strand_B_tapered_dust_preparation.tex, review_B_tapered_dust_preparation.tex;
% the corresponding deliveries/solve_b/*_scope.tex files.
% Their named matching algebra/source-error claims, distinct preparations,
% adjoint allocation and unsolved coupled-history scope are retained above.
% The new exact native vertex spectrum and center/exterior trace-class
% construction were saved before reading the RED continuation. Its full
% current source/SOURCES was then read: the independent star reduction,
% Gaussian derivative bound and global relative-mass contraction agree
% with the local-center plus exterior route above. The RED's precise
% angular obstruction was independently checked and incorporated. These
% new results remain [R]; no external matched isolated collapse or
% exact unprojected radial state has been supplied by this construction.
% The complete current declared_h2_continuum.tex and all SOURCES were
% read after this independent continuation. Its finite-parent center
% contraction, normalized lapse and conserved-norm residence argument
% agree after sigma_here=-sigma_H2. Its open-parent fractional positive-
% band preparation and source projection are distinct from the localized
% infinite-parent preparation here; both retain the angular and geometry-
% scale conditions. This same-lane comparison is [R], not an external [X].
